% Word count aim: 4000-4500
% current: 3165
\chapter{Background}
\label{cha:backg}


\section{Computational Problems}

\subsection{Problem Formulations}
% Define problem instances and problem classes.
% Distinguish between decision, search, and optimization problems.
% Clarify deterministic vs. stochastic problem formulations.

A \textit{computational problem} is a formal specification of a task defined over a set of
\textit{problem instances}. Each instance encodes the input data relevant to a particular
occurrence of the problem. The \textit{problem class} describes the collection of all
such problem instances together with the corresponding output requirements.~\cite{HOOS200513}
 
Problems are broadly categorized by the nature of their required output. 
\begin{itemize}
    \item A \textit{decision problem} demands a binary answer, whether a given instance satisfies some property. 
    \item A \textit{search problem} asks for the construction of an explicit solution
    satisfying a specified property, such as a feasible assignment or a path in a graph.
    \item An \textit{optimization problem} further requires that the returned solution be optimal (or near-optimal) with respect to some objective function.~\cite{HOOS200513}
\end{itemize}

A further distinction concerns the role of randomness in the problem formulation. In a
\textit{deterministic} formulation, both the instance and every algorithm operating on it are fully specified without any probabilistic component. In a \textit{stochastic}
formulation, the instance itself, the objective evaluations, or the algorithm may involve randomness. Noisy objective functions, where repeated evaluation of the same candidate
solution may return different values, are a canonical example and arise naturally in
simulation-based and real-world settings.

\subsection{Algorithms}
\subsubsection{Algorithm Properties}
\paragraph{Deterministic and Randomized Algorithms. }
% Definitions and differences.

\paragraph{Adaptive vs. Non-Adaptive Algorithms}
% Sequential decision processes.
% Role of feedback.

\subsubsection{Oracle-Based Computation}
% Definition of oracle access.
% Query-response interaction model.

An \textit{oracle} is an abstract computational device that answers queries about a
problem instance. In optimization, the most common oracle is the \textit{function
evaluation oracle}, which, given a query point $x \in \mathcal{X}$, returns the
objective value $f(x)$ (or a noisy estimate thereof). The algorithm interacts with the
oracle in a \textit{query-response} model: it submits a query and receives a response
before deciding on the next query. The total number of such interactions constitutes the
\textit{query complexity} of the algorithm, formalized in
Section~\ref{subsubsec:infocomplexity}~\cite{nocedal2006numerical, droste2006blackbox}.



\subsection{Search Space}
\subsubsection{Structure of the Search Space}
% Finite vs. infinite spaces.
% Discrete vs. continuous domains.
% Combinatorial vs. geometric domains.

The \textit{search space} $\mathcal{X}$ is the set of all candidate solutions considered
by an algorithm. Search spaces may be \textit{finite} or \textit{infinite}, and
\textit{discrete} or \textit{continuous}. 
 
Discrete combinatorial search spaces, in particular the Boolean hypercube $\{0,1\}^n$, are central to the theoretical analysis of evolutionary algorithms (\acrshort{ea}) and other randomized search
heuristics~\cite{droste2002analysis,HOOS200561} used in theoretical \acrshort{bbo} analysis. 
Continuous spaces $\mathbb{R}^d$ arise in numerical optimization~\cite{nocedal2006numerical}, and mixed spaces containing both discrete and continuous variables occur in practice but are considerably harder to
analyze~\cite{audet2017derivative}.
 
The structure of the search space, and in particular how neighboring solutions relate to
one another, fundamentally influences the difficulty of optimization. This is encoded by
a \textit{neighborhood function} $\mathcal{N} : \mathcal{X} \to 2^{\mathcal{X}}$,
assigning to each candidate solution the set of solutions reachable in a single step!CITE: Hoos and Stützle - 2005 - INTRODUCTION!.
For the Boolean hypercube, the standard neighborhood consists of all strings at Hamming
distance one !CITE: a resource using hamming distance!.

\begin{remark}[Problem domain of interest]
This project is concerned with \emph{Pseudo-Boolean, Combinatorial Optimization problems}, with solution space $\mathcal{X} \subseteq \{0,1\}^n$. A candidate solution can be represented as a \emph{bitstring}: a sequence of alphabets $a \in \{0,1\}$. The objective functions considered for such problems are of the form $f:\mathcal{X} \to \mathbb{R}_0^+$, where $\mathbb{R}_0^+$ denotes the non-negative reals. Without loss of generality, an objective function with co-domain $\mathcal{T} \subseteq \mathbb{R}$ can be mapped to this convention via the affine shift
\[
    g: \mathbb{R} \to \mathbb{R}_0^+, \qquad g(x) = x - \min_{y \in \mathcal{X}} f(y),
\]
which is well-defined whenever $f$ attains its minimum on $\mathcal{X}$, and preserves the ordering of all objective values by construction.
\end{remark}


\subsubsection{Constraints and Feasibility}
% Feasible region definition.
% Constrained vs. unconstrained problems.

In a \textit{constrained} optimization problem, the \textit{feasible region}
$\mathcal{F} \subseteq \mathcal{X}$ is defined by one or more constraints that a valid
solution must satisfy~\cite{nocedal2006numerical, bertsekas2009convex}. 
An \textit{unconstrained} problem places no restriction beyond
membership in $\mathcal{X}$. Constraints may be \textit{hard} (a solution violating them
is inadmissible) or \textit{soft} (violations are penalized in the objective). The
boundary of the feasible region can significantly complicate optimization, particularly
for local search methods that may generate infeasible intermediate solutions.

\subsection{Objective Functions and Landscape Properties}

The geometric structure of the objective function determines both which algorithms are
applicable and what convergence guarantees are attainable. A central organizing
principle is that access to higher-order analytical information about $f$, i.e., gradients, 
curvature, or structural properties such as convexity, generally enables stronger
theoretical guarantees and more efficient algorithms; restricting this access,
as in the black-box setting, correspondingly limits what can be guaranteed.
 
\subsubsection{Zeroth-, First-, and Second-Order Methods}
% Keep existing content; add forward reference to §2.2 (BBO = zeroth-order world)
 
Optimization algorithms can be classified by the order of analytical information they
require from $f$. This \textit{information hierarchy} has three tiers.
 
\textit{Zeroth-order} methods use only evaluations of $f(x)$, with no gradient or
curvature information. The algorithm observes only the output of the function evaluation
oracle at each queried point and must infer search directions from comparisons of
objective values alone. This exactly describes the black-box model formalized in
Section~\ref{sec:bbo}: when the internal structure of $f$ is inaccessible or
analytically intractable, zeroth-order methods are the only option. Their weaker
information access is the fundamental reason black-box optimization is harder in
general than white-box optimization, and it is the regime in which all algorithms
studied in this project operate.
 
\textit{First-order} methods additionally access the gradient $\nabla f(x) \in
\mathbb{R}^d$, which encodes the direction of steepest ascent at $x$. For a
differentiable function $f : \mathbb{R}^d \to \mathbb{R}$, gradient-based methods
exploit this to define descent directions and form the backbone of classical continuous
optimization~\cite{nocedal2006numerical}. When $f$ is convex but non-differentiable,
the gradient generalizes to the \textit{subdifferential} $\partial f(x)$, the set of
all \textit{subgradients} $g$ satisfying
\[
  f(y) \geq f(x) + g^\top(y - x) \quad \forall\, y \in \mathbb{R}^d.
\]
Subgradient methods replace the gradient with any element of $\partial f(x)$, recovering
convergence under convexity alone, albeit at rates slower than those attainable with
smooth gradients~\cite{bertsekas2009convex}.
 
\textit{Second-order} methods further access the Hessian $\nabla^2 f(x) \in
\mathbb{R}^{d \times d}$, which captures local curvature. Methods that utilize or
approximate the Hessian, such as Newton and quasi-Newton schemes, can achieve
superlinear or quadratic convergence near a solution~\cite{nocedal2006numerical}, at the
cost of maintaining a $d \times d$ matrix.
 
Together, these three tiers define distinct levels of oracle access, and their
availability directly governs the choice and efficiency of applicable algorithms. The
move from second-order to zeroth-order corresponds to progressively weaker assumptions
about what the oracle reveals, culminating in the black-box model where only function
values are observable.
 
\subsubsection{Landscape Properties}
 
The following geometric properties of $f$ characterize which tier of analytical
information is both available and useful.
 
\paragraph{Convexity.}
A function $f : \mathbb{R}^d \to \mathbb{R}$ is \textit{convex} if for all
$x, y \in \mathbb{R}^d$ and $\lambda \in [0,1]$~\cite{bertsekas2009convex},
\[
  f\!\left(\lambda x + (1-\lambda)y\right) \leq \lambda f(x) + (1-\lambda)f(y).
\]
Convexity guarantees that every local minimum is global, rendering the landscape
tractable for first-order methods. \textit{Strong convexity} additionally requires
\[
  f(y) \geq f(x) + \nabla f(x)^\top(y - x) + \tfrac{\mu}{2}\|y - x\|^2
\]
for some $\mu > 0$, implying a unique global minimum and linear convergence rates for
gradient descent~!CITE: Nesterov 2004!.
 
\paragraph{Smoothness and Lipschitz Continuity.}
% TODO
 
\paragraph{Ruggedness and Multimodality.}
% - Multiple local optima, deceptive gradients, epistatic interactions
% - Motivates heuristic / stochastic methods: no first-order structure to exploit
% - Forward reference to NK-landscapes as a tunable model of ruggedness

\subsection{Complexity Measures}
\subsubsection{Classical Complexity: Time and Space}
% Standard asymptotic measures.

\subsubsection{Information Complexity: Query and Oracle Complexity}
\label{subsubsec:infocomplexity}
% Number of oracle calls as complexity measure.

\subsubsection{Expected vs. Worst-Case Analysis}
% Runtime notions in randomized algorithms.


\subsection{Probabilistic Tools for Algorithm Analysis}
\subsubsection{Expectation and Variance}
\subsubsection{Concentration Inequalities}
\paragraph{Markov's Inequality. }
\paragraph{Chebyshev's Inequality. }
\paragraph{Chernoff Bounds. }

\newpage

\section{Black-Box Optimization}
\label{sec:bbo}

\subsection{The Black-Box Model}
% Formal oracle model.
% Zeroth-order information.
% Information revealed per query.

In the \textit{black-box optimization} model, the algorithm has no access to the
analytical form of the objective function. It interacts with the objective exclusively
through an oracle that, upon receiving a query $x \in \mathcal{X}$, returns $f(x)$ (or
a noisy estimate thereof)~\cite{droste2006blackbox,doerr2020theory,lengler2019drift}. 
This constitutes \textit{zeroth-order} information: no gradient or higher-order information is available.
The black-box model is appropriate whenever the objective arises from a simulation,
physical experiment, or any process whose internal structure is inaccessible or too
complex to exploit analytically.


\subsubsection{Variants of Black-Box Models}
% Unrestricted model.
% Ranking-based model.
% Memory-restricted model.


\subsection{White-Box vs. Black-Box Optimization}
% - What structural knowledge white-box methods exploit
% - Gradient availability as the dividing line
% - Computational trade-offs: information richness vs. access cost

\subsection{Black-Box Complexity}
\label{subsec:bb-complexity}
% - Definition: minimum number of queries any algorithm requires in the worst case
%   to solve a given problem class  [droste2006blackbox]
% - Yao's Minimax Principle: lower bounds via distributions over instances
% - Upper bounds (algorithmic) and lower bounds (information-theoretic)
% - Comparison-based lower bounds (sorting analogy)
% - What these results say about the hardness of combinatorial BBO

\subsection{Purpose of Randomized Search Heuristics}
% Explain why, given the BBO setting, one turns to general-purpose
% stochastic methods rather than problem-specific exact algorithms.
%
% - No-Free-Lunch Theorems: statement and implications for algorithm comparison
%   [wolpert1997no or similar]
%   No universally superior algorithm exists across all problem instances;
%   the choice of algorithm embeds implicit assumptions about problem structure.
% - Practical consequences:
%     When the function structure is unknown, heuristic methods are natural
%     When exact methods are computationally infeasible (NP-hard instances),
%     anytime heuristics offer the best available solution within a budget
%     Stochasticity enables escape from local optima that deterministic
%     local search cannot avoid
% - Taxonomy of derivative-free / black-box methods:
%     Exact DFO (Nelder-Mead, pattern search, model-based DFO)
%     Bayesian Optimization (Gaussian process surrogate + acquisition)
%     Evolutionary / population-based heuristics
%     Single-solution local search heuristics
%     Simulated Annealing and other acceptance-criterion methods
% - Scope of this project: randomized single- and population-based heuristics
%   on combinatorial / pseudo-Boolean domains
 
\subsubsection{No-Free-Lunch Theorems}
 
\subsubsection{Taxonomy of Derivative-Free Methods}
% Brief map of the full space; help the reader locate RSHs within it


\subsection{Randomized Search Heuristic Algorithms}


\subsubsection[\texorpdfstring{Evolutionary Algorithms (EA)}{Evolutionary Algorithms (EA)}]{Evolutionary Algorithms (\acrshort{ea})}



\subsubsection{Simulated Annealing}


\paragraph{Cooling Schedules}



\subsubsection[\texorpdfstring{Randomized Local Search (RLS)}{Randomized Local Search (RLS)}]{\acrfull{rls}}




\paragraph[\texorpdfstring{Stochastic Hill-Climber - Greedy RLS}{Stochastic Hill-Climber - Greedy RLS}]{Stochastic Hill-Climber - Greedy \acrshort{rls}}



\paragraph[\texorpdfstring{$\epsilon$-greedy RLS}{epsilon-greedy RLS}]{$\epsilon-greedy$ \acrshort{rls}}



\paragraph[\texorpdfstring{Decaying $\epsilon-greedy$ RLS}{Decaying epsilon-greedy RLS}]{Decaying $\epsilon-greedy$ \acrshort{rls}}



\subsection{Few Pseudo-Boolean Problems}

The \textit{pseudo-Boolean} setting considers functions $f : \{0,1\}^n \to \mathbb{R}$
defined on binary strings of length $n$. Several benchmark functions are widely used to
study and compare the behavior of randomized search heuristics, each encoding a
different type of difficulty.~\cite{lengler2019drift,droste2006blackbox}

\subsubsection{OneMax}

The \textsc{OneMax} function counts the number of one-bits in a bit string:
\[
  \textsc{OneMax}(x) = \sum_{i=1}^{n} x_i, \qquad x \in \{0,1\}^n.
\]
The unique global maximum is $x^{*} = 1^n$ with $\textsc{OneMax}(1^n) = n$.
\textsc{OneMax} is the canonical baseline benchmark: the landscape is unimodal and every
improving move increases the function value by exactly one, providing a clear gradient
for local search methods. The expected runtime of the $(1+1)$~\acrfull{ea} on \textsc{OneMax} is
$\Theta(n \log n)$. !CITE: Droste, Jansen, Wegener (2002)!

\subsubsection{LeadingOnes}

The \textsc{LeadingOnes} function counts the length of the maximal leading prefix of
ones:
\[
  \textsc{LeadingOnes}(x) = \sum_{i=1}^{n} \prod_{j=1}^{i} x_j.
\]
Equivalently, $\textsc{LeadingOnes}(x)$ equals the index of the first zero bit minus
one, or $n$ if no zero exists. The unique global optimum is again $1^n$. Because
progress requires the leading bit to be set before any subsequent bit can contribute,
\textsc{LeadingOnes} is significantly harder than \textsc{OneMax}: the expected runtime
of the $(1+1)$~\acrshort{ea} is $\Theta(n^2)$. !CITE: Droste, Jansen, Wegener (2002)!


\subsubsection{Jump}

The \textsc{Jump} function with parameter $k \in \{1, \ldots, n\}$ coincides with
\textsc{OneMax} on bit strings with at most $n - k$ ones and on the global optimum
$1^n$, but introduces a region of reduced fitness for strings with between $n-k+1$ and
$n-1$ ones. This creates a ``gap'' of width $k$ that algorithms relying on single-bit
improvements cannot cross: the set of strings with exactly $n-k$ ones constitutes a
local optimum of standard local search, and reaching the global optimum $1^n$ requires
flipping all $k$ remaining zero bits simultaneously. For the $(1+1)$~\acrshort{ea} with standard
bit mutation at rate $1/n$, the probability of a specific $k$-bit flip is
$\Theta(n^{-k})$, giving an expected runtime of $\Theta(n^k)$ on $\textsc{Jump}_k$.

Formally, \textsc{Jump}$_{k,n} \colon \{0,1\}^n \to \mathbb{R}$ is defined as
\[
  \textsc{Jump}_{k,n}(x) \;=\;
  \begin{cases}
    k + \displaystyle\sum_{i=1}^{n} x_i & \text{if } \displaystyle\sum_{i=1}^{n} x_i \leq n - k
    \text{ or } \displaystyle\sum_{i=1}^{n} x_i = n, \\[8pt]
    n - \displaystyle\sum_{i=1}^{n} x_i & \text{otherwise.}
  \end{cases}
\]
Droste, Jansen, and Wegener~\cite{droste2002analysis} proved that the expected running
time of the $(1+1)$~\acrshort{ea} on $\textsc{Jump}_{k,n}$ is $\Theta(n^k + n\log n)$ for
$k \in \{1, 2, \ldots, n\}$, which simplifies to $\Theta(n^k)$ for $k \geq 2$.
\subsubsection{BinaryValue[BinVal]}
The \textsc{BinaryValue} (\textsc{BinVal}) function interprets a bit string as a binary
integer by weighting bits exponentially by position:
\[
  \textsc{BinVal}(x) = \sum_{i=1}^{n} 2^{i-1}\, x_i.
\]
The global optimum is $1^n$ with value $2^n - 1$. The exponential weighting means that
the most significant bit contributes more than the combined value of all less significant
bits, making \textsc{BinVal} highly structured.

\subsubsection{Trap}
\subsubsection{Plateau}
\subsubsection{Hierarchical If-Only-If}


\subsubsection{MaxSAT}
\textit{Maximum Satisfiability} (MaxSAT) is the optimization variant of the Boolean
satisfiability problem (SAT). Given a propositional formula in conjunctive normal form (CNF) !TODO: Cite or explain CNF!
with $m$ clauses over $n$ Boolean variables, the goal is to find a truth assignment that
maximizes the number of satisfied clauses. MaxSAT is NP-hard in general~\cite{hutchison_search_2004}
and serves as an important benchmark for combinatorial optimization algorithms.

\subsubsection{Tunable Landscapes: NK-Landscapes}

\textit{NK-landscapes}~\cite{kauffman1987towards} are a parametric family of pseudo-Boolean
functions controlled by two integers: $N$ (the number of variables) and $K$ (the number
of epistatic interactions per variable, $0 \leq K \leq N-1$). The fitness of a string
$x \in \{0,1\}^N$ is defined as
\[
  f(x) = \frac{1}{N} \sum_{i=1}^{N} f_i\!\left(x_i,\, x_{\pi(i,1)},
      \ldots,\, x_{\pi(i,K)}\right),
\]
where each local contribution $f_i : \{0,1\}^{K+1} \to [0,1]$ is drawn independently
and uniformly at random, and $\pi(i, \cdot)$ denotes the $K$ epistatic neighbours of
variable $i$.

When $K = 0$ the contributions are independent, producing a smooth, additively separable
landscape analogous to \textsc{OneMax}. As $K$ increases toward $N-1$, epistatic
interactions intensify and the landscape becomes increasingly rugged, with the number of
local optima growing in the limit. NK-landscapes thus provide a tunable model of
ruggedness and are widely used to study algorithmic scalability and robustness across a
spectrum of problem difficulties.

\subsection{Traditional Performance Measures in BBO}
% - Expected runtime (hitting time to reach the optimum or a quality threshold)
%   as the standard measure in EA theory  [droste2002analysis; doerr2020theory]
% - Fixed-budget analysis: quality of best-found solution after T queries
%   [Jansen & Zarges 2014 or similar]
% - Fixed-target analysis: minimum budget to reach a target quality with
%   probability >= 1-delta  [!CITE!]
% - Limitations of these measures:
%     Binary in nature (found / not found)
%     Ignore the trajectory of quality improvement
%     Do not penalise "wasteful" queries on already-known-poor regions
% - This gap motivates regret as a more informative measure


\newpage
\section{Sequential and Online Optimization Frameworks}

\subsection{Exploration-Exploitation Trade-off}
%   Gathering information (exploration) vs. using current best
%   knowledge (exploitation)
% - Why pure exploitation fails (local optima / suboptimal arms)
% - Why pure exploration fails (never converges to good solutions)
% - This tension is shared with BBO

\subsection{Multi-Armed Bandits}
% - Formal stochastic bandit model: K arms, reward distributions, T rounds
% - Objective: maximize cumulative reward (equivalently, minimize regret)
% - Stochastic MAB: UCB1, Thompson Sampling — brief intuition  [!CITE!]
% - Contextual bandits: features available at each round  [!CITE!]
% - Pure exploration / best-arm identification: goal is simple regret
%   at the end of the horizon, not cumulative regret  [bubeck2011pure]
% - Connection to BBO: each candidate solution is an "arm"; the oracle
%   is the reward; the algorithm decides which arm to pull

\subsection{Online Optimization}
% Sequential decision-making setup.
% Adversarial vs. stochastic settings.



% TODO: Confirm
\subsection{Reinforcement Learning}
% Value functions and policies.
% Policy optimization perspective.

\newpage
\section{Computational Regret}

\subsection{Definition of Regret}
% Formal mathematical definition.
% Regret vs. loss.
Let $f : \mathcal{X} \to \mathbb{R}$ be an objective function with global optimum value
$f^{*} = \max_{x \in \mathcal{X}} f(x)$. An algorithm that queries points
$x_1, x_2, \ldots, x_T$ incurs \textit{regret} by querying suboptimal solutions. Regret
quantifies the gap between the algorithm's performance and that of an oracle always
querying the optimum~\cite{cesa2006prediction}.
!TODO: add more context with reference to MAB \& rewards!

\subsection{Types of Regret}

\subsubsection{Instantaneous Regret}
% Per-round performance gap.
The \textit{instantaneous regret} at round $t$ is
\[
  \rho_t = f^{*} - f(x_t),
\]
capturing the suboptimality of the single query made at step $t$.
A related variant tracks the \textit{incumbent} (best-so-far) solution rather than
the current query:
\[
  \rho^{\mathrm{inc}}_t = f^{*} - f(\hat{x}_t),
  \quad \hat{x}_t = \operatorname{arg\,max}_{s \leq t} f(x_s).
\]
Because $\hat{x}_t$ is non-decreasing in quality, $\rho^{\mathrm{inc}}_t$ is
monotonically non-increasing over $t$, and is equivalent to simple regret evaluated at time $t$(~\ref{subsubsec:sr}).

\subsubsection{Cumulative Regret}
\label{subsubsec:cr}
% Definition, interpretation, and relationship to IR.
Two cumulative regret variants arise naturally and serve different purposes.
 
The \textit{\acrfull{cr}} after $T$ rounds is~\cite{cesa2006prediction}
\[
  R_T = \sum_{t=1}^{T} \rho_t = \sum_{t=1}^{T} \bigl(f^{*} - f(x_t)\bigr).
\]
This is the standard formulation from the online learning literature: it sums the
instantaneous suboptimality of each individual query, regardless of what the algorithm
has previously observed. Minimizing $R_T$ corresponds to performing well
\textit{throughout} the optimization process, penalizing every suboptimal query equally.
 
The \textit{incumbent cumulative regret} instead tracks the best-so-far process
\[
  B(t) \;:=\; \max_{s \leq t} f(x_s),
\]
yielding
\[
  R^{\mathrm{inc}}_T = \sum_{t=1}^{T} \rho^{\mathrm{inc}}_t = \sum_{t=1}^{T}
  \bigl(f^{*} - B(t)\bigr).
\]
Since $B(t)$ is non-decreasing in $t$, the per-step contributions
$\rho^{\mathrm{inc}}_t$ are non-increasing, and $R^{\mathrm{inc}}_T \leq R_T$ for all
$T$. The incumbent formulation is the appropriate notion when the practitioner retains
the best solution found rather than only the most recent query, which is the standard
convention for the randomized search heuristics studied in this project. All regret
analyses in Chapter~\ref{cha:exper} use $R^{\mathrm{inc}}_T$ unless otherwise stated.
 
In continuous time, the incumbent cumulative regret is approximated via a left-hold
(piecewise-constant) integral
\[
  R^{\mathrm{inc}}_T \approx \sum_{t=1}^{T} \rho^{\mathrm{inc}}_{t-1}\,\Delta t_{\,t},
  \quad \Delta t_{\,t} = t - t_{t-1},
\]
which accounts for unequal spacing between evaluations when function evaluation time is
non-uniform. In all experiments reported here, evaluations are unit-cost and time is
measured in discrete query counts, so $\Delta t_t = 1$ throughout and this reduces to
the standard discrete sum.

\subsubsection{Simple Regret}
% Final recommendation quality.
\label{subsubsec:sr}
The \textit{\acrfull{sr}} after $T$ rounds is~\cite{bubeck2011pure}
\[
  r_T = f^{*} - f(\hat{x}_T),
\]
where $\hat{x}_T = \operatorname{arg\,max}_{t \leq T} f(x_t)$ is the best solution
encountered. Simple regret measures the quality of the \textit{final output} and is the
appropriate notion when the goal is to return a good solution rather than to perform well
during the search. Note that $r_T = \rho^{\mathrm{inc}}_T$: simple regret is precisely the incumbent instantaneous regret at the final round.

When comparing algorithms across multiple independent runs, the standard scalar summary
is the \textit{expected simple regret}
\[
  \mathbb{E}[r_T] = \frac{1}{N}\sum_{i=1}^{N} r_T^{(i)},
\]
where $r_T^{(i)}$ denotes the simple regret of the $i$-th run.

\subsubsection{Normalized Regret}
% Scale-free comparison across problem instances.
To enable comparison across problem instances with different objective scales, simple
regret can be normalized by the range of attainable fitness values:
\[
  \tilde{r}_T = \frac{f^{*} - f(\hat{x}_T)}{f^{*} - f_{\mathrm{worst}}} \in [0, 1],
\]
where $f_{\mathrm{worst}}$ is the worst observed or theoretical fitness value. A value of
$0$ indicates an optimal solution was found; a value of $1$ indicates no improvement over
the worst case. When $f^{*} = f_{\mathrm{worst}}$ (a degenerate constant landscape), the
problem is trivially solved and normalized regret is defined as $0$.



% Expectation over prior distributions.

\subsection{Regret in Black-Box Optimization}
% Regret under query constraints.
% Regret under noisy oracle access (Aside).
% Regret as fine-grained runtime measure.
Regret provides a fine-grained complexity measure in the black-box setting: rather than
simply asking whether an algorithm finds the optimum within $T$ queries, regret tracks
how suboptimal the queries are throughout the process. This is particularly relevant
when every evaluation is costly and performing well before the budget is exhausted
matters in practice.

\subsection{Regret Bounds}
% Upper bounds.
% Lower bounds.
% Minimax regret.
% Instance-dependent bounds.

\subsection{Connections to Runtime Analysis}
% Relationship between regret and optimization time.
% Trade-offs between exploration and convergence.
% Regret as a refined complexity notion.

\newpage

